\chapter{Fazit \& Ausblick}
\label{kap:Fazit und Ausblick}
Im Rahmen dieser Studienarbeit konnte erfolgreich die methodische und technische Grundlage für die Transformation des vorhandenen Robotersystems - bestehend aus dem EX RV-M1 und einem Vision-System (2D \& 3D-ToF) - hin zu einem limitiert kollaborativen und dynamischen System gelegt werden.
\\
\\Zunächst wurden durch eine IST-Analyse die konkreten Limitationen der synchronen Softwarearchitektur, insbesondere die blockierenden Wartezeiten durch die Verwendung von Thread.Sleep()-Befehlen, quantitativ aufgezeigt. Auf Basis dieser Erkenntnisse wurde eine asynchrone Befehlsstruktur konzipiert und implementiert. Diese verwendet moderne .NET-Prinzipien (TPL und TAP) und eine prädiktive Zeitberechnung mittels inverser Kinematik, um die Laufzeit des Systems an die physisch auszuführenden Bewegung anzupassen. Diese Umstellung bietet weitere Optimierungspotenziale und Testfälle, welche in nachfolgenden Arbeiten umgesetzt werden können. Dazu zählen:
\begin{itemize}
	\item Weitere Optimierung der Laufzeit durch Messung der real benötigten Zeit für das Senden von Befehlen sowie Öffnen/Schließen des Greifers
	\item Aufbrechen der MP-Befehle über Trajektoriensegmentierung um ein schnelleres Stoppen der Pipeline nach jedem gefahrenen Segment zu realisieren (vgl. Konzept in Studienarbeit T3100 Adrian Sommer \& Luis Thiel)
	\item Grenzfallbetrachtung: Tests mit verschiedenen Objekten und Orientierungen - Zuverlässigkeit und Reproduzierbarkeit evaluieren
\end{itemize}
Des Weiteren konnte eine aktive Detektion von menschlichen Händen mithilfe der 2D-Kamera umgesetzt werden. Hierzu kommt die Machine-Learning-Methode \textit{Haar Cascade Classifier} zum Einsatz. Dabei handelt es sich um einen vortrainierten Datensatz, welcher es ermöglicht, einen fest definierten Arbeitsbereich zu überwachen und diesen nach Merkmalen zu filtern, welche einer Hand entsprechen. Die finale Implementation wurde erfolgreich mittels einer Testreihe validiert, in welcher auch ein Ignorieren der für den Roboter relevanten Objekte gezeigt wurde.\\
\\Aufbauend auf dieser Studienarbeit ergeben sich für die Zukunft Weiterentwicklungspotenziale wie z.B. die Untersuchung von KI-basierten Safety-Ansätzen oder die Implemenation eines externen Sensor/CPU-System um einen hardwareseitigen Not-Halt in die Steuerung zu implementieren.
